\section{infinite products}
%%%
\index{infinite products}

Just as we developed the theory for infinite sums, we could develop the theory
for infinite products by setting
\[
  \prod_{k=0}^{+\infty} a_k = \lim_{n\to +\infty} \prod_{k=0}^n a_k
  = \lim_{n\to +\infty} a_0 \cdot a_1 \cdots a_n.
\]

We will always assume $a_k>0$, otherwise the sign of the product would be
difficult to define. Then, by using the logarithm (which transforms products
into sums), we can relate infinite products to series:
\[
  \prod_{k=0}^{+\infty} a_k = \lim_{n\to +\infty} e^{\sum_{k=0}^n \ln a_k}.
\]

Observe that if the series of logarithms diverges to $-\infty$, the infinite
product has a limit of $0$. Having required that the terms $a_k$ are all
positive, the product can never be less than zero. To maintain the analogy with
series, we will say that the infinite product converges if the limit of the
partial products is finite and positive. We will say it diverges if the limit is
$+\infty$ or $0$.

Thus, we can say that the infinite product converges if and only if the series
of logarithms converges.

We then observe that a necessary condition for an infinite product $\prod a_k$
to be convergent is that $\ln a_k\to 0$, i.e., $a_k \to 1$. In this case, since
\[
  \ln a_k = \ln (1+(a_k-1)) \sim a_k - 1
\]
it is observed that if $a_k\to 1$ and $a_k\ge 1$, the infinite product $\prod
a_k$ converges if and only if the series $\sum (a_k-1)$ converges.

\begin{example}[sum of the reciprocals of primes]
\index{primes!sum of the reciprocals}%
\index{sum!of the reciprocals of primes}%
We can use infinite products to demonstrate that the sum of the reciprocals of
prime numbers is divergent. Let $p_k$ be the sequence of prime numbers ($p_1=2$,
$p_2=3$, $p_3=5$, $\dots$ we are assuming that prime numbers are infinite). Then
we want to demonstrate that
\begin{equation}\label{eq:489467523}
  \sum_{k=1}^{+\infty} \frac{1}{p_k} = +\infty.
\end{equation}

This result is significant in number theory as it tells us that $p_k$ cannot go
to infinity like a power $k^\alpha$ with $\alpha>1$ since the series $\sum
1/k^\alpha$ is convergent.

We now show that~\eqref{eq:489467523} holds. Note that for every $n\in \NN$, the
term $\frac 1 n$ can be decomposed as the product of powers of the reciprocals
of prime numbers. The idea is to consider the product:
\begin{align*}
  \prod_{k=1}^{N} \sum_{j=0}^{M} \frac{1}{p_k^j}
  &= 
  (1 + \frac 1 2 + \frac 1 {2^2} + \dots + \frac{1}{2^M}) \cdot
  (1 + \frac 1 3 + \frac 1 {3^2} + \dots + \frac{1}{3^M}) \cdot \\
  &\quad \cdot(1 + \frac 1 5 + \frac 1 {5^2} + \dots + \frac{1}{5^M}) \cdots
  (1 + \frac 1 {p_N} + \frac 1 {p_N^2} + \dots + \frac{1}{p_N^M})
\end{align*}
observing that in expanding this product, we obtain as addends the reciprocals
of the products of all powers of the prime numbers $p_1,\dots, p_N$ with
exponents not exceeding $M$. Thus, if $n < p_{N+1}$ and $n \le 2^M$, then $n$ is
one of such products and thus $\frac 1 n$ appears in the above expansion.
Therefore, we can assert that
\begin{equation}\label{eq:8834884}
  \prod_{k=1}^{+\infty} \frac{1}{1-\frac{1}{p_k}}
  = \prod_{k=1}^{+\infty} \sum_{j=0}^{+\infty} \frac{1}{p_k^j}
  \ge \prod_{k=1}^{N} \sum_{j=0}^{M} \frac{1}{p_k^j}
  \ge \sum_{n=1}^{K-1} \frac 1 n
\end{equation}
where $K$ is the smallest between $p_{N+1}$ and $2^M$. Since $K$ can be made
arbitrarily large by choosing $N$ and $M$ sufficiently large, and since $\sum
\frac 1 n =+\infty$, it means that the product on the left side of
\eqref{eq:8834884} must itself be divergent. We then recall that this product
has the same nature as the following series
\[
  \sum_{k=1}^{+\infty} \enclose{\frac{1}{1-\frac 1 {p_k}}-1}
  = \sum_{k=1}^{+\infty} \frac{\frac 1 {p_k}}{1-\frac{1}{p_k}}
\]
but since $1/p_k\to 0$, we have
\[
  \frac{\frac 1 {p_k}}{1-\frac{1}{p_k}}
  \sim \frac 1 {p_k}
\]
and thus, by the asymptotic comparison test, the previous series has the same
nature as the series
\[
\sum_{k=1}^{+\infty} \frac{1}{p_k}
\]
which is therefore divergent.
\end{example}

%%%